%% Un livre et un ballon posés sur une table, en perspective. Repère : y vertical, le plateau
%% est le plan (x, z). Les deux pièces ont PLUSIEURS mobilités, et pas le même nombre : le livre
%% ① glisse dans les deux directions du plateau et pivote à plat (appui plan, 3), le ballon ②
%% peut en plus rouler dans toutes les directions (ponctuelle, 5). Le nombre de cases à cocher
%% ne se devine pas : il se lit sur le contact.
\documentclass[tikz,border=4pt]{standalone}
\usepackage{liaisons}
%% _objets-common.tex — outils communs aux figures « objet du quotidien avec son repère ».
%%
%% Ces figures servent des exercices où l'axe ne doit PAS être donné dans l'énoncé : l'élève lit
%% l'orientation du repère sur le dessin. D'où deux exigences :
%%   - le repère porte ses TROIS axes, dans une orientation qui change d'une figure à l'autre
%%     (si le vertical était toujours z, « vertical → z » redeviendrait un réflexe) ;
%%   - tout est en \large : à 390 px de large sur un téléphone, \small est illisible.
%%
%% Le \repere de liaisons.sty ne convient pas ici : deux axes seulement, et en \small.
%%
%% RÈGLE DE COULEUR, et elle est stricte : la couleur dit **mobile ou fixe** par rapport au
%% référentiel de l'exercice, jamais « la pièce dont on parle ». Le bâti — sol, route, mur, meuble,
%% tapis, plateau, support — est en solideB (bleu) ou en gris ; TOUT ce qui peut bouger par rapport
%% à lui est en solideA (rouge), même quand plusieurs pièces bougent. Ce sont les pastilles ① ② qui
%% désignent la pièce interrogée, c'est leur seule raison d'être. Colorer « la pièce ① » en rouge et
%% « la pièce ② » en bleu induirait en erreur une fois sur deux l'élève qui a appris que le rouge
%% bouge — le cas s'est produit sur le vélo, dont le cadre était bleu alors qu'il roule.

%% \repereXYZ[longueur]{point}{angle/label, angle/label, angle/label}
%% Trois flèches depuis un même point, chacune à l'angle voulu. L'axe « qui sort de la feuille »
%% se dessine comme les autres, en diagonale, à la manière des perspectives du cours.
\newcommand\repereXYZ[3][0.95]{%
  \foreach \ang/\lab in {#3} {%
    \draw[-{Stealth[length=6pt]}, line width=0.8pt] #2 -- ++(\ang:#1)
      node[pos=1, anchor={\ang+180}, inner sep=3pt, font=\large] {$\vec{\lab}$};%
  }}

%% \numero{point du repère}{point visé sur le dessin}{n} — pastille noire reliée à la pièce.
\newcommand\numero[3]{%
  \draw[line width=0.5pt, black!55] #1 -- #2;
  \node[circle, fill=black, text=white, inner sep=2pt, font=\large\bfseries] at #1 {#3};}

%% Épaisseurs : le dessin doit tenir à la taille d'un téléphone, donc peu de traits et des traits gras.
\tikzset{
  objet/piece/.style={line width=1.6pt, line cap=round, line join=round, draw=solideA},
  objet/bati/.style={line width=1.6pt, line cap=round, line join=round, draw=solideB},
  objet/fin/.style={line width=0.7pt, line cap=round, draw=black!55},
}

\begin{document}
\begin{tikzpicture}[x=1cm,y=1cm]

\coordinate (prof) at (28:1.05);   % calc n'accepte pas une macro ici

% --- le plateau de la table : le bâti
\draw[objet/bati, fill=solideB!5]
  (0,0) -- (5.4,0) -- ($(5.4,0)+(prof)$) -- ($(0,0)+(prof)$) -- cycle;
\draw[objet/bati] (0.25,0) -- (0.25,-1.15);
\draw[objet/bati] (5.15,0) -- (5.15,-1.15);

% --- le livre, posé à plat : il glisse et pivote sur le plateau
% dessiné en volume : à plat il ne se distinguait pas du plateau
\coordinate (li) at (0.85,0.12);
\draw[objet/piece, fill=solideA!8]
  (li) -- ++(1.55,0) -- ++(28:0.78) -- ++(-1.55,0) -- cycle;
\draw[objet/piece, fill=solideA!14]
  (li) -- ++(1.55,0) -- ++(0,0.34) -- ++(-1.55,0) -- cycle;
\draw[objet/piece] ($(li)+(1.55,0.34)$) -- ++(28:0.78);
% la tranche : trois traits fins, pour qu'on lise « livre » et non « boîte »
\foreach \d in {0.10,0.17,0.24} { \draw[objet/fin] ($(li)+(0.06,\d)$) -- ++(1.43,0); }

% --- le ballon, posé dessus : il roule en plus de glisser
\coordinate (ba) at (4.05,0.85);
\draw[objet/piece, fill=white] (ba) circle (0.55);
\draw[objet/fin] ($(ba)+(-0.55,0)$) arc (180:360:0.55 and 0.2);

% --- repère : y vertical, le plateau est le plan (x, z)
\repereXYZ{(6.35,0.35)}{90/y, 0/x, 28/z}

\numero{(0.55,2.05)}{($(li)+(0.75,0.5)$)}{1}
\numero{(5.55,2.25)}{($(ba)+(0.35,0.42)$)}{2}

\end{tikzpicture}
\end{document}
